\phantomsection
\chapter*{Dadan dan Ujang}
\addcontentsline{toc}{section}{Dadan dan Ujang --- Reyhan Eka}
\penulis{Reyhan Eka}

\awalcerpen{D}{adan dan Ujang} adalah teman sekelas yang sudah lama bersahabat. Mereka selalu bersama ketika pergi dan pulang sekolah. Rumah mereka juga tidak terlalu jauh. Mereka sering berangkat bersama pada pagi hari. Dadan dikenal sebagai anak yang baik, ramah, dan suka menolong teman. Ujang adalah anak yang rajin dan bertanggung jawab terhadap tugas-tugas sekolahnya. Walaupun mereka punya sifat yang berbeda, keduanya tetap saling menghargai dan tidak pernah mempermasalahkan perbedaan itu.

Suatu pagi, seperti biasa Dadan menunggu Ujang di depan rumahnya. Setelah beberapa menit, Ujang datang. Wajahnya terlihat sedikit terburu-buru.

``Dan, ayo cepat. Nanti kita terlambat,'' kata Ujang.

``Iya, tunggu sebentar. Kita pergi sekarang,'' jawab Dadan.

Mereka lalu berjalan ke sekolah sambil bicara tentang pelajaran yang akan dipelajari hari itu. Pagi itu cuacanya cukup cerah. Banyak siswa lain juga tampak berjalan ke sekolah. Ada yang bersama teman-temannya. Ada juga yang diantar orang tua.

Ketika hampir sampai di sekolah, tiba-tiba Ujang berhenti di pinggir jalan. Ia membuka tasnya, lalu memeriksa semua buku yang dibawanya. Wajahnya mulai panik.

``Kenapa, Jang?'' tanya Dadan.

``Buku tugasku tidak ada, Dan. Padahal hari ini harus dikumpulkan,'' jawab Ujang dengan cemas.

Dadan berusaha menenangkan Ujang. Ia lalu membantu memeriksa isi tas Ujang. Mereka mengeluarkan buku satu per satu. Namun, buku tugas yang dicari tetap tidak ditemukan.

``Coba kamu ingat lagi. Tadi malam kamu mengerjakan tugasnya di mana?'' tanya Dadan.

``Di rumah. Setelah selesai, aku masukkan ke buku. Tapi aku sudah lupa setelah itu,'' jawab Ujang.

Ujang semakin khawatir karena tugas itu harus dikumpulkan kepada guru. Ia takut dimarahi karena tidak membawa tugas.

``Jangan khawatir, Jang. Kita cari bersama. Mungkin bukunya jatuh di jalan,'' kata Dadan.

Akhirnya mereka memutuskan untuk kembali menyusuri jalan yang tadi mereka lewati. Mereka berjalan perlahan. Mereka melihat ke pinggir jalan. Dadan memperhatikan setiap tempat yang mungkin menjadi tempat jatuhnya buku.

Mereka pertama-tama melihat di dekat warung yang mereka lewati. Namun, buku itu tidak ada. Mereka lalu melihat di dekat pagar dan tempat duduk di pinggir jalan. Hasilnya tetap sama.

``Belum ketemu juga,'' kata Ujang dengan wajah kecewa.

``Jangan menyerah dulu. Kita cari lagi. Ingat, tadi kita berhenti di beberapa tempat,'' jawab Dadan.

Mereka terus berjalan. Setelah beberapa saat, Dadan melihat sesuatu yang kelihatan seperti buku di bawah sebuah pohon.

``Jang, coba lihat itu!'' seru Dadan. Ia menunjuk ke arah pohon.

Ujang langsung mendekatinya. Ternyata benar. Buku yang ada di bawah pohon itu adalah buku tugas miliknya. Mungkin buku itu jatuh saat mereka berjalan.

``Alhamdulillah, ketemu!'' kata Ujang dengan wajah senang.

Ujang mengambil bukunya dan memeriksa keadaan buku itu. Untungnya, buku itu masih baik, hanya sedikit kotor karena terkena tanah.

``Untung kamu melihatnya, Dan. Kalau tidak, mungkin aku tidak akan menemukan buku ini,'' kata Ujang.

``Tidak apa-apa. Yang penting bukunya sudah ketemu. Sekarang kita cepat ke sekolah supaya tidak terlambat,'' jawab Dadan.

Mereka lalu segera pergi ke sekolah. Saat tiba di sekolah, bel masuk hampir berbunyi. Mereka langsung masuk ke kelas dan duduk di tempat masing-masing.

Tidak lama kemudian, guru masuk ke kelas. Guru itu meminta semua siswa mengumpulkan tugas yang sudah diberikan sebelumnya. Ujang langsung mengeluarkan buku tugasnya. Ia lalu menyerahkannya kepada guru.

``Ini tugas saya, Bu,'' kata Ujang.

Guru menerima buku itu lalu mengangguk.

``Baik, Ujang. Terima kasih sudah mengumpulkan tepat waktu,'' jawab guru.

Ujang merasa lega. Ia lalu melihat ke arah Dadan dan tersenyum. Ia merasa sangat bersyukur karena punya teman seperti Dadan yang mau membantunya saat ia sedang mengalami kesulitan.

Saat jam istirahat, Ujang mendekati Dadan.

``Dan, terima kasih ya sudah membantu aku tadi pagi,'' kata Ujang.

``Iya, sama-sama. Kita kan teman. Kalau teman sedang kesusahan, kita harus bantu,'' jawab Dadan.

Ujang lalu berkata, ``Aku jadi sadar kalau punya teman itu bukan hanya untuk bermain atau bercanda. Teman juga harus saling membantu.''

Dadan mengangguk. ``Benar. Kita juga tidak tahu kapan kita akan butuh bantuan orang lain.''

Sejak kejadian itu, persahabatan Dadan dan Ujang jadi makin dekat. Mereka lebih sering saling membantu dalam belajar atau kegiatan sehari-hari. Jika Dadan kesulitan memahami pelajaran, Ujang akan membantunya. Begitu juga jika Ujang butuh bantuan, Dadan selalu berusaha membantu selama ia bisa.

Beberapa hari kemudian, Dadan sendiri kesulitan saat mengerjakan tugas kelompok. Ia lupa membawa beberapa bahan yang dibutuhkan untuk tugas itu. Ujang tahu hal itu lalu langsung menawarkan bantuan.

``Dan, kamu bisa pakai punyaku. Nanti kita kerjakan bersama,'' kata Ujang.

Dadan tersenyum. ``Terima kasih, Jang.''

``Sama-sama. Dulu kamu juga sudah menolong aku,'' jawab Ujang.

Dadan merasa senang. Kebaikan yang ia lakukan membuat Ujang belajar melakukan hal yang sama kepada orang lain. Mereka lalu mengerjakan tugas bersama sampai selesai.

Dari kejadian-kejadian itu, Dadan dan Ujang belajar bahwa persahabatan bukan hanya soal sering bersama, bermain bersama, atau tertawa bersama. Persahabatan juga soal saling membantu, saling menghargai, dan tetap ada di samping teman saat menghadapi masalah.

Sejak saat itu, mereka mencoba menjadi teman yang lebih baik. Mereka tidak ingin hanya membantu saat keadaan mudah. Mereka juga mau membantu ketika situasi sedang sulit. Menurut mereka, persahabatan akan lebih berarti jika ada rasa peduli dan saling membantu.

\jedahias
